\scanpage{70}
% archive sheet number: 69
% @ L1–7
B) Recollements \inserted{ne} donnant \uline{pas} lieu à une application naturelle $\mathscr X\to X$, i.e. l'isom $\sigma:\alpha\to\beta$ est \uline{opposé} à \unsure{celle|celui} obtenu par transp. par., ou par composition de celui-ci avec antipodisme -- ou encore, l'isom de transition $\underline{\omega}(D)\to\underline{\omega}(D')$ est \deleted{opposé à} celui du transport parallèle. Deux cas à envisager encore

% @ L8–14
1) Cas $\alpha=a'$. C'est le cas où on répère une op. de générisation \inserted{sur $D'$} par l'arc qui fait face à l'arc de spécialisation correspondant sur $D$ -- c'était un cas que pour une raison ou une autre j'avais écarté (c'était vraiment au bonheur-la-chance, rétrospectivement !).

% @ L15–end to R2
2) Cas $\alpha=a$ -- C'est le cas où on repère une op. de générisation par l'arc \inserted{de gén.} déduit de l'arc de spécialisation correspondant par transport parallèle. Je me suis amplement convaincu qu'il est impossible de ``suivre'' une p.v. en lui appliquant les itérés de l'opération \deleted{\unsure{\ensuremath{\rho}|\ensuremath{\beta}}} de ``tourner autour d'un trou (ou sommet \ldots).

% @ R3–14
Je crois me rappeler maintenant que j'avais écarté aussi le cas 1), car il donnait lieu aussi à des circuits qui avaient l'air idiots -- notamment en partant d'une p.v. proche d'une supersingulière, $D_i$, on s'écarte de $D_i$ au lieu de tourner autour comme on devrait normalement.

% @ R foot, cancelled lines; last line continues
\deleted{Là je viens de revérifier, cet ennui se présente aussi bien pour B 1) que pour B 2), et \unsure{ça|\ensuremath{=}} pour \deleted{\unsure{B}} A 1), et même, je viens de voir,} pour A 2 -- visiblement, dans ce dernier
